% Esercitazione TEEA 3MAT — Argomento 2: partitori, potenza ed energia
% Per rigenerare il PDF (dalla cartella strumenti, due volte):
%   pdflatex -output-directory=../../../materiali -jobname=esercitazione-partitori-potenza-energia-3MAT esercitazione.tex
% ---------------------------------------------------------------------------
%  Stile comune di esercitazione e soluzioni (pdfLaTeX).
%  Viene incluso da esercitazione.tex e soluzioni.tex con \input.
% ---------------------------------------------------------------------------
\documentclass[11pt,a4paper]{article}
\usepackage[a4paper, top=14mm, bottom=18mm, left=16mm, right=16mm, footskip=9mm]{geometry}
\usepackage[utf8]{inputenc}
\usepackage[T1]{fontenc}
\usepackage{helvet}
\renewcommand{\familydefault}{\sfdefault}
\usepackage{amsmath, amssymb, graphicx, xcolor, fancyhdr, lastpage, tabularx, booktabs, array, enumitem}
\usepackage[most]{tcolorbox}
\usepackage[siunitx, american resistors, american currents]{circuitikz}
\ctikzset{bipoles/thickness=1, current/distance=0.35}
\usetikzlibrary{calc, arrows.meta}

\definecolor{titolo}{HTML}{2A7FA8}
\definecolor{grigio}{HTML}{E9ECEF}
\definecolor{tenue}{HTML}{555555}
\setlength{\parindent}{0pt}
\setlength{\parskip}{3pt}

\newcommand{\U}[1]{\,\mathrm{#1}}
\newcommand{\virg}{{,}}
\newcommand{\riga}[1][28mm]{\underline{\hspace{#1}}}

% intestazione come nelle esercitazioni della scuola
\newcommand{\intestazione}[2]{%
  \begin{minipage}[c]{20mm}\includegraphics[width=18mm]{logo-isis.png}\end{minipage}%
  \begin{minipage}[c]{\dimexpr\linewidth-20mm}
    {\fontsize{20}{24}\selectfont\bfseries\color{titolo} #1}\par\vspace{1mm}
    {\footnotesize\bfseries Indirizzo: IP14 -- MANUTENZIONE E ASSISTENZA TECNICA \quad -- \quad Classe 3MAT}
  \end{minipage}\par\vspace{2mm}
  {\color{titolo}\rule{\linewidth}{0.8pt}}\par\vspace{3mm}
  \begin{tcolorbox}[colback=grigio, colframe=grigio, arc=0pt, left=6pt, top=4pt, bottom=4pt]
    {\Large\bfseries #2}
  \end{tcolorbox}}

\newcommand{\nota}[1]{\begin{tcolorbox}[colback=grigio!45, colframe=grigio, arc=0pt, left=6pt, top=3pt, bottom=3pt]\small #1\end{tcolorbox}}
\newcommand{\tipo}[1]{\par\vspace{2mm}{\large\color{titolo} #1}\par\vspace{-1mm}{\color{titolo!60}\rule{\linewidth}{0.4pt}}\par\vspace{1mm}}
\newcommand{\esercizio}[2]{\par\vspace{3mm}\textbf{Esercizio #1.} \textbf{#2}\par}

% alimentazione: barra con etichetta
\newcommand{\rail}[2]{\draw[line width=1.2pt] ($#1+(-0.6,0)$) -- ($#1+(0.6,0)$); \node[above] at #1 {#2};}

\pagestyle{fancy}
\fancyhf{}
\renewcommand{\headrulewidth}{0pt}
\newcommand{\piede}[1]{\fancyfoot[L]{\scriptsize I.S.I.S. Sacile-Brugnera | sede: ``Piero Della Valentina'' | #1}
  \fancyfoot[R]{\scriptsize Pagina \thepage\ di \pageref{LastPage}}}

% Schemi degli esercizi (usati sia nell'esercitazione sia nelle soluzioni)
\newcommand{\partitoreV}[3]{% #1 = V, #2 = R1, #3 = R2
\begin{circuitikz}[scale=0.8, transform shape, line width=0.6pt]
  \rail{(0,4.4)}{$V = #1$}
  \draw (0,4.4) -- (0,4.2) to[R, l_={$R_1 = #2$}] (0,2.2) to[short, *-] (0,2.2)
        to[R, l_={$R_2 = #3$}] (0,0.2) -- (0,0) node[sground]{};
  \draw (0,2.2) to[short, -o] (1.4,2.2) node[right]{X};
\end{circuitikz}}

\newcommand{\partitoreI}[3]{% #1 = I, #2 = R1, #3 = R2
\begin{circuitikz}[scale=0.8, transform shape, line width=0.6pt]
  \draw (0,4.4) node[ocirc]{} to[short, i>^={$I = #1$}] (0,3.2);
  \draw (-1.6,3.2) -- (1.6,3.2); \draw (0,3.2) node[circ]{};
  \draw (-1.6,3.2) to[short, i>^={$I_1$}] (-1.6,2.6) to[R, l_={$R_1 = #2$}] (-1.6,0.6) -- (-1.6,0.3) -- (1.6,0.3);
  \draw (1.6,3.2) to[short, i>_={$I_2$}] (1.6,2.6) to[R, l^={$R_2 = #3$}] (1.6,0.6) -- (1.6,0.3);
  \draw (0,0.3) node[circ]{} -- (0,0.1) node[sground]{};
\end{circuitikz}}

\newcommand{\partitorePot}{%
\begin{circuitikz}[scale=0.8, transform shape, line width=0.6pt]
  \rail{(0,4.4)}{$V = 10\,\mathrm{V}$}
  \draw (0,4.4) -- (0,4.2) to[R, l_={$R_1 = 1\,\mathrm{k\Omega}$}] (0,2.2) to[short, *-] (0,2.2)
        to[vR, l_={$R_2 = 0 \div 4{,}7\,\mathrm{k\Omega}$}] (0,0.2) -- (0,0) node[sground]{};
  \draw (0,2.2) to[short, -o] (1.4,2.2) node[right]{X};
\end{circuitikz}}

\newcommand{\resistenzaSola}[2]{% #1 = V, #2 = R
\begin{circuitikz}[scale=0.8, transform shape, line width=0.6pt]
  \rail{(0,3.4)}{$V = #1$}
  \draw (0,3.4) -- (0,3.2) to[R, l_={$R = #2$}, i>^={$I$}] (0,0.8) -- (0,0.6) node[sground]{};
\end{circuitikz}}

\piede{Esercitazione TEEA 3MAT}

\begin{document}

\intestazione{ESERCITAZIONE DI TEEA}{Partitori, potenza ed energia}

\nota{Esercizi di allenamento su partitore di tensione, partitore di corrente, potenziometro, potenza ed energia.
Scrivi sempre la formula, sostituisci i valori con le unità di misura e controlla il risultato con le verifiche
($V_1 + V_X = V$, $I_1 + I_2 = I$). L'energia va espressa in Wh o kWh.}

\vspace{2mm}
\textbf{COGNOME E NOME:} \riga[95mm] \hfill \textbf{CLASSE:} \riga[22mm]

\tipo{Partitore di tensione}

\esercizio{1}{Nel partitore di tensione in figura calcola la corrente $I$, la tensione $V_1$ su $R_1$, la tensione
$V_X$ nel punto X e la potenza dissipata da ciascuna resistenza e dal circuito.}
\begin{minipage}[t]{0.38\linewidth}\vspace{0pt}\centering
\partitoreV{+24\,\mathrm{V}}{3{,}3\,\mathrm{k\Omega}}{4{,}7\,\mathrm{k\Omega}}
\end{minipage}\hfill
\begin{minipage}[t]{0.58\linewidth}\vspace{2mm}
\renewcommand{\arraystretch}{2.1}
\begin{tabular}{@{}l l@{}}
$I =$ \riga & $V_1 =$ \riga\\
$V_X =$ \riga & $P_1 =$ \riga\\
$P_2 =$ \riga & $P_{tot} =$ \riga\\
\end{tabular}
\end{minipage}
\vspace{18mm}

\esercizio{2}{Si vuole ottenere $V_X = 5\U{V}$ da un'alimentazione di $15\U{V}$ con un partitore in cui $R_2 = 2{,}2\U{k\Omega}$.
Calcola il valore di $R_1$, poi scegli il valore commerciale più vicino della serie E12
(\dots\ 3{,}3 -- 3{,}9 -- 4{,}7 -- 5{,}6 -- 6{,}8\ \dots\ k$\Omega$) e ricalcola la tensione $V_X$ che si ottiene.}
\begin{minipage}[t]{0.38\linewidth}\vspace{0pt}\centering
\partitoreV{+15\,\mathrm{V}}{\,?}{2{,}2\,\mathrm{k\Omega}}
\end{minipage}\hfill
\begin{minipage}[t]{0.58\linewidth}\vspace{2mm}
\renewcommand{\arraystretch}{2.1}
\begin{tabular}{@{}l@{}}
$R_1$ calcolata $=$ \riga[35mm]\\
$R_1$ commerciale $=$ \riga[35mm]\\
$V_X$ con $R_1$ commerciale $=$ \riga[30mm]\\
\end{tabular}
\end{minipage}
\vspace{14mm}

\newpage
\tipo{Partitore di corrente}

\esercizio{3}{Nel partitore di corrente in figura calcola le correnti $I_1$ e $I_2$, la resistenza equivalente
$R_{eq}$ e la tensione $V$ ai capi del parallelo.}
\begin{minipage}[t]{0.46\linewidth}\vspace{0pt}\centering
\partitoreI{30\,\mathrm{mA}}{2\,\mathrm{k\Omega}}{3\,\mathrm{k\Omega}}
\end{minipage}\hfill
\begin{minipage}[t]{0.50\linewidth}\vspace{2mm}
\renewcommand{\arraystretch}{2.1}
\begin{tabular}{@{}l@{}}
$I_1 =$ \riga \\ $I_2 =$ \riga \\ $R_{eq} =$ \riga \\ $V =$ \riga\\
\end{tabular}
\end{minipage}
\vspace{12mm}

\esercizio{4}{Nel partitore di corrente in figura calcola $I_1$, $I_2$ e la tensione $V$ ai capi del parallelo.
Calcola poi la potenza dissipata da $R_1$ e da $R_2$: per $R_1$ basta un resistore da $1/4\U{W}$,
tenendo un margine di almeno il doppio?}
\begin{minipage}[t]{0.46\linewidth}\vspace{0pt}\centering
\partitoreI{50\,\mathrm{mA}}{100\,\Omega}{400\,\Omega}
\end{minipage}\hfill
\begin{minipage}[t]{0.50\linewidth}\vspace{2mm}
\renewcommand{\arraystretch}{2.1}
\begin{tabular}{@{}l l@{}}
$I_1 =$ \riga & $I_2 =$ \riga\\
$V =$ \riga & \\
$P_1 =$ \riga & $P_2 =$ \riga\\
\multicolumn{2}{@{}l@{}}{$1/4\U{W}$ basta? \riga[40mm]}\\
\end{tabular}
\end{minipage}
\vspace{12mm}

\newpage
\tipo{Potenziometro}

\esercizio{5}{Nel circuito in figura $R_2$ è un potenziometro che si regola da $0$ a $4{,}7\U{k\Omega}$.
a) Calcola la tensione $V_X$ minima e massima e la corrente massima nel circuito.
b) Se invece volessi regolare $V_X$ da circa $2\U{V}$ fino a $10\U{V}$, dove metteresti il potenziometro?
Con $R_2 = 1\U{k\Omega}$ fissa, quanto dovrebbe valere al massimo il potenziometro?}
\begin{minipage}[t]{0.40\linewidth}\vspace{0pt}\centering
\partitorePot
\end{minipage}\hfill
\begin{minipage}[t]{0.56\linewidth}\vspace{2mm}
\renewcommand{\arraystretch}{2.1}
\begin{tabular}{@{}l l@{}}
a) $V_{X\,min} =$ \riga[22mm] & $V_{X\,max} =$ \riga[22mm]\\
\phantom{a)} $I_{max} =$ \riga[22mm] & \\
b) posizione: \riga[30mm] & $R_{1\,max} =$ \riga[22mm]\\
\end{tabular}
\end{minipage}
\vspace{14mm}

\tipo{Potenza ed energia}

\esercizio{6}{Un resistore da $470\U{\Omega}$ è collegato a $12\U{V}$. Calcola la corrente e la potenza dissipata.
Va bene un resistore da $1/4\U{W}$? Quale potenza nominale sceglieresti (1/4 W, 1/2 W o 1 W), tenendo un margine di almeno il doppio?}
\begin{minipage}[t]{0.30\linewidth}\vspace{0pt}\centering
\resistenzaSola{12\,\mathrm{V}}{470\,\Omega}
\end{minipage}\hfill
\begin{minipage}[t]{0.66\linewidth}\vspace{2mm}
\renewcommand{\arraystretch}{2.1}
\begin{tabular}{@{}l l@{}}
$I =$ \riga & $P =$ \riga\\
$1/4\U{W}$ va bene? \riga[25mm] & potenza scelta: \riga[22mm]\\
\end{tabular}
\end{minipage}
\vspace{10mm}

\newpage
\esercizio{7}{In una resistenza da $100\U{\Omega}$ scorre una corrente di $0{,}2\U{A}$ per $5$ ore.
Calcola la tensione ai capi della resistenza, la potenza e l'energia consumata in Wh e in kWh.}
\vspace{2mm}
\renewcommand{\arraystretch}{2.1}
\begin{tabular}{@{}l l l@{}}
$V =$ \riga & $P =$ \riga & $E =$ \riga[22mm] Wh $=$ \riga[22mm] kWh\\
\end{tabular}
\vspace{30mm}

\esercizio{8}{Una stufetta elettrica da $1500\U{W}$ funziona a $230\U{V}$ per $4$ ore al giorno, per $30$ giorni.
Calcola la corrente assorbita, la resistenza della stufetta, l'energia consumata nel mese in Wh e in kWh
e il costo, con un prezzo di $0{,}25$ euro al kWh.}
\vspace{2mm}
\renewcommand{\arraystretch}{2.1}
\begin{tabular}{@{}l l@{}}
$I =$ \riga & $R =$ \riga\\
$E =$ \riga[30mm] Wh $=$ \riga[22mm] kWh & costo $=$ \riga[22mm] euro\\
\end{tabular}
\vspace{30mm}

\end{document}
